\documentclass[10pt,a4paper]{amsart}
\usepackage{amsmath,amssymb,amsthm,graphicx}
% Retained minimal preamble for the Stacks Project chapter sites-cohomology.tex (Cohomology on Sites).
% Verbatim excerpt of the chapter preamble: Xy-pic setup (native preamble lines 14-18) and the native
% theorem declarations with the mathematical macros (native preamble lines 154-217). Copied mechanically;
% nothing added, removed or redefined. The wrapper class and the amsmath/amssymb/amsthm/graphicx base
% packages are supplied by the assembler, so the class/externaldocument part of the native preamble is omitted.
% For commutative diagrams we use Xy-pic
\usepackage[all]{xy}

% We use 2cell for 2-commutative diagrams.
\xyoption{2cell}
% Theorem environments.
%
\theoremstyle{plain}
\newtheorem{theorem}[subsection]{Theorem}
\newtheorem{proposition}[subsection]{Proposition}
\newtheorem{lemma}[subsection]{Lemma}

\theoremstyle{definition}
\newtheorem{definition}[subsection]{Definition}
\newtheorem{example}[subsection]{Example}
\newtheorem{exercise}[subsection]{Exercise}
\newtheorem{situation}[subsection]{Situation}

\theoremstyle{remark}
\newtheorem{remark}[subsection]{Remark}
\newtheorem{remarks}[subsection]{Remarks}

\numberwithin{equation}{subsection}

% Macros
%
\def\lim{\mathop{\mathrm{lim}}\nolimits}
\def\colim{\mathop{\mathrm{colim}}\nolimits}
\def\Spec{\mathop{\mathrm{Spec}}}
\def\Hom{\mathop{\mathrm{Hom}}\nolimits}
\def\Ext{\mathop{\mathrm{Ext}}\nolimits}
\def\SheafHom{\mathop{\mathcal{H}\!\mathit{om}}\nolimits}
\def\SheafExt{\mathop{\mathcal{E}\!\mathit{xt}}\nolimits}
\def\Sch{\mathit{Sch}}
\def\Mor{\mathop{\mathrm{Mor}}\nolimits}
\def\Ob{\mathop{\mathrm{Ob}}\nolimits}
\def\Sh{\mathop{\mathit{Sh}}\nolimits}
\def\NL{\mathop{N\!L}\nolimits}
\def\CH{\mathop{\mathrm{CH}}\nolimits}
\def\proetale{{pro\text{-}\acute{e}tale}}
\def\etale{{\acute{e}tale}}
\def\QCoh{\mathit{QCoh}}
\def\Ker{\mathop{\mathrm{Ker}}}
\def\Im{\mathop{\mathrm{Im}}}
\def\Coker{\mathop{\mathrm{Coker}}}
\def\Coim{\mathop{\mathrm{Coim}}}

% Boxtimes
%
\DeclareMathSymbol{\boxtimes}{\mathbin}{AMSa}{"02}

%
% Macros for moduli stacks/spaces
%
\def\QCohstack{\mathcal{QC}\!\mathit{oh}}
\def\Cohstack{\mathcal{C}\!\mathit{oh}}
\def\Spacesstack{\mathcal{S}\!\mathit{paces}}
\def\Quotfunctor{\mathrm{Quot}}
\def\Hilbfunctor{\mathrm{Hilb}}
\def\Curvesstack{\mathcal{C}\!\mathit{urves}}
\def\Polarizedstack{\mathcal{P}\!\mathit{olarized}}
\def\Complexesstack{\mathcal{C}\!\mathit{omplexes}}
% \Pic is the operator that assigns to X its picard group, usage \Pic(X)
% \Picardstack_{X/B} denotes the Picard stack of X over B
% \Picardfunctor_{X/B} denotes the Picard functor of X over B
\def\Pic{\mathop{\mathrm{Pic}}\nolimits}
\def\Picardstack{\mathcal{P}\!\mathit{ic}}
\def\Picardfunctor{\mathrm{Pic}}
\def\Deformationcategory{\mathcal{D}\!\mathit{ef}}

% --- appended 2026-10-12 (Chief): the project own definitions of the reference/slogan/history
% display environments, copied mechanically from the shared Stacks preamble lines 8-12
% (\usepackage{verbatim} plus the three \newenvironment lines, which the chapter preamble keeps
% verbatim from the project preamble). Nothing added, removed or redefined: these environments
% carry no mathematical content and the project itself typesets them as comments. Without these
% three lines a counted statement whose own text uses them fails with Environment slogan undefined.
\usepackage{verbatim}
\newenvironment{reference}{\comment}{\endcomment}
\newenvironment{slogan}{\comment}{\endcomment}
\newenvironment{history}{\comment}{\endcomment}
\usepackage{cleveref}
\crefname{originaltheorem}{original theorem}{original theorems}
\Crefname{originaltheorem}{Original theorem}{Original theorems}
\makeatletter
\newlabel{injectives-theorem-K-injective-embedding-grothendieck}{{079P}{source}{}{sourcepointer-1}{}}
\newlabel{injectives-theorem-K-injective-embedding-grothendieck@cref}{{[originaltheorem][0][]079P}{[1][0][]source}{}{}{}}
\makeatother
\crefname{originallemma}{original lemma}{original lemmas}
\Crefname{originallemma}{Original lemma}{Original lemmas}
\makeatletter
\newlabel{lemma-restriction-along-monomorphism-surjective}{{093X}{source}{}{sourcepointer-2}{}}
\newlabel{lemma-restriction-along-monomorphism-surjective@cref}{{[originallemma][0][]093X}{[1][0][]source}{}{}{}}
\makeatother
\crefname{originallemma}{original lemma}{original lemmas}
\Crefname{originallemma}{Original lemma}{Original lemmas}
\makeatletter
\newlabel{derived-lemma-adjoint-preserve-K-injectives}{{08BJ}{source}{}{sourcepointer-3}{}}
\newlabel{derived-lemma-adjoint-preserve-K-injectives@cref}{{[originallemma][0][]08BJ}{[1][0][]source}{}{}{}}
\makeatother
\crefname{originallemma}{original lemma}{original lemmas}
\Crefname{originallemma}{Original lemma}{Original lemmas}
\makeatletter
\newlabel{injectives-lemma-derived-products}{{07D9}{source}{}{sourcepointer-4}{}}
\newlabel{injectives-lemma-derived-products@cref}{{[originallemma][0][]07D9}{[1][0][]source}{}{}{}}
\makeatother
\title{The derived limit of an inverse system agrees with the homotopy limit of its terms}
\author{The Stacks Project Authors}
\date{stacks-project:sites-cohomology.tex:local-snapshot-20261011; GNU FDL 1.2 or later}
\begin{document}
\maketitle
\noindent\textit{Editorial note.} All delivered mathematical text is the author's own, copied from the pinned Stacks Project chapter file sites-cohomology.tex; no statement or proof step was rewritten, reordered or removed.\par\medskip
\noindent\textit{Editorial note.} The proof is the author's complete argument for the identification, including the surjectivity of the transition maps and the comparison with derived products; the closing reference to the definition of homotopy limits is an identification, not an omission.\par\medskip
\noindent\textit{Editorial note.} Every labelled reference occurring inside the counted spans is carried in assembly.external\_references as a non-load-bearing navigation pointer (same-chapter ones at their definition site, cross-chapter ones at the invocation inside the counted span), so no load-bearing mathematics is replaced by a stub. No bibliography citation occurs inside the counted spans.\par\medskip
\noindent\textit{Editorial note.} The checked chapter list for same-name collisions is the set of stacks chapters staged in this repository: chow, cohomology, divisors, etale-cohomology, flat, more-algebra, more-morphisms, morphisms, spaces-more-morphisms, varieties together with this card's chapter sites-cohomology.\par\medskip
\noindent\textit{Editorial note.} The preamble is the chapter's own retained minimal preamble (sources/stacks-sites-cohomology/retained-preamble.tex), a verbatim excerpt of the native preamble: the Xy-pic setup and the native theorem declarations with the mathematical macros.\par\medskip
\noindent\textit{Allowed standard prerequisites (editorial).}\begin{itemize}
\item The derived limit as a derived pushforward. Same chapter, the description of R lim as Rg\_* for the morphism of topoi g associated to the projection from the product site. Supplies the functor R lim being computed.
\item Existence of K-injective representatives with injective terms. Injectives, Theorem injectives-theorem-K-injective-embedding-grothendieck. Supplies the representative with which R lim is computed termwise.
\item Products in the derived category. Injectives, Lemma injectives-lemma-derived-products together with Derived Categories, Lemma derived-lemma-adjoint-preserve-K-injectives. Identifies the termwise products with derived-category products and keeps the truncations K-injective.
\end{itemize}\medskip
\ifdefined\hypertarget\hypertarget{sourcepointer-1}{}\fi
\noindent\textit{Source reference.} Original theorem 079P: sources/stacks-sites-cohomology/pinned-original-source.tex, source lines 5082--5082. Navigation pointer invoked at proof line 5082; the counted span contains the pointer only and the cited result is not reproduced inside it, so the pointer is registered as a non-load-bearing navigation stub in the established form for Stacks material. No bibliography citation is involved.\par\medskip
\ifdefined\hypertarget\hypertarget{sourcepointer-2}{}\fi
\noindent\textit{Source reference.} Original lemma 093X: sources/stacks-sites-cohomology/pinned-original-source.tex, source lines 1962--1962. Navigation pointer invoked at proof line 5096; the counted span contains the pointer only and the cited result is not reproduced inside it, so the pointer is registered as a non-load-bearing navigation stub in the established form for Stacks material. No bibliography citation is involved.\par\medskip
\ifdefined\hypertarget\hypertarget{sourcepointer-3}{}\fi
\noindent\textit{Source reference.} Original lemma 08BJ: sources/stacks-sites-cohomology/pinned-original-source.tex, source lines 5106--5106. Navigation pointer invoked at proof line 5106; the counted span contains the pointer only and the cited result is not reproduced inside it, so the pointer is registered as a non-load-bearing navigation stub in the established form for Stacks material. No bibliography citation is involved.\par\medskip
\ifdefined\hypertarget\hypertarget{sourcepointer-4}{}\fi
\noindent\textit{Source reference.} Original lemma 07D9: sources/stacks-sites-cohomology/pinned-original-source.tex, source lines 5119--5119. Navigation pointer invoked at proof line 5119; the counted span contains the pointer only and the cited result is not reproduced inside it, so the pointer is registered as a non-load-bearing navigation stub in the established form for Stacks material. No bibliography citation is involved.\par\medskip
\medskip\noindent
On the other hand, the continuous and cocontinuous functors
$\mathcal{C} \to \mathcal{C} \times \mathbf{N}$,
$U \mapsto (U, n)$ define morphisms of topoi
$i_n : \Sh(\mathcal{C}) \to \Sh(\mathcal{C} \times \mathbf{N})$.
Of course $i_n^{-1}$ is the functor which picks the $n$th term of
the inverse system. Thus there are transformations of functors
$i_{n + 1}^{-1} \to i_n^{-1}$. Hence given
$K \in D(\mathcal{C} \times \mathbf{N})$ we get
$K_n = i_n^{-1}K \in D(\mathcal{C})$ and maps $K_{n + 1} \to K_n$.

\begin{lemma}
\label{lemma-derived-limit-is-ok}
Let $\mathcal{C}$ be a site. Let $K$ be an object of
$D(\mathcal{C} \times \mathbf{N})$. Set $K_n = i_n^{-1}K$ as above.
Then
$$
R\lim K \cong R\lim K_n
$$
in $D(\mathcal{C})$.
\end{lemma}

\begin{proof}
To calculate $R\lim$ on an object $K$ of $D(\mathcal{C} \times \mathbf{N})$
we choose a K-injective representative $\mathcal{I}^\bullet$ whose terms are
injective objects of $\textit{Ab}(\mathcal{C} \times \mathbf{N})$, see
Injectives, Theorem
\ref{injectives-theorem-K-injective-embedding-grothendieck}.
We may and do think of $\mathcal{I}^\bullet$ as an inverse system of
complexes $(\mathcal{I}_n^\bullet)$ and then we see that
$$
R\lim K = \lim \mathcal{I}_n^\bullet
$$
where the right hand side is the termwise inverse limit.

\medskip\noindent
Let $\mathcal{J} = (\mathcal{J}_n)$ be an injective object of
$\textit{Ab}(\mathcal{C} \times \mathbf{N})$. The morphisms
$(U, n) \to (U, n + 1)$ are monomorphisms of
$\mathcal{C} \times \mathbf{N}$, hence
$\mathcal{J}(U, n + 1) \to \mathcal{J}(U, n)$ is surjective
(Lemma \ref{lemma-restriction-along-monomorphism-surjective}).
It follows that $\mathcal{J}_{n + 1} \to \mathcal{J}_n$ is
surjective as a map of presheaves.

\medskip\noindent
Note that the functor $i_n^{-1}$ has an exact left adjoint $i_{n, !}$.
Namely, $i_{n, !}\mathcal{F}$ is the inverse system
$\ldots 0 \to 0 \to \mathcal{F} \to \ldots \to \mathcal{F}$.
Thus the complexes $i_n^{-1}\mathcal{I}^\bullet = \mathcal{I}_n^\bullet$
are K-injective by
Derived Categories, Lemma \ref{derived-lemma-adjoint-preserve-K-injectives}.

\medskip\noindent
Because we chose our K-injective complex to have injective terms
we conclude that
$$
0 \to  \lim \mathcal{I}_n^\bullet \to \prod \mathcal{I}_n^\bullet
\to \prod \mathcal{I}_n^\bullet \to 0
$$
is a short exact sequence of complexes of abelian sheaves as it
is a short exact sequence of complexes of abelian presheaves.
Moreover, the products in the middle and the right represent
the products in $D(\mathcal{C})$, see
Injectives, Lemma \ref{injectives-lemma-derived-products} and its
proof (this is where we use that $\mathcal{I}_n^\bullet$ is K-injective).
Thus $R\lim K$ is a homotopy limit of the inverse system $(K_n)$
by definition of homotopy limits in triangulated categories.
\end{proof}

\end{document}
